problem:  
**Problem (Existence and Universality of the Second Spectral Term for Hörmander Systems).**  
Let \( (M, \mathcal H, g_{\mathcal H}) \) be a compact smooth sub‑Riemannian manifold where \( \mathcal H\subset TM \) is a bracket‑generating distribution satisfying Hörmander’s condition with fixed step \(r\). Let  
\[
\Delta_{\mathcal H} = -\sum_{i=1}^k X_i^*X_i
\]
be the associated sub‑Laplacian acting on scalar functions, and let \(N(\lambda)\) denote the counting function of its eigenvalues below \(\lambda\).

It is known that the leading asymptotic (“Weyl law”) holds:
\[
N(\lambda) \sim C_{\mathcal H}\,\lambda^{Q/2}, \qquad \lambda\to\infty,
\]
where \(Q\) is the homogeneous dimension and \(C_{\mathcal H}\) depends on the osculating nilpotent group at each point.

**Main open problem:**  
1. (*Existence*) Does there exist, for general bracket‑generating \(\mathcal H\), a second‑order term \(c_2(\mathcal H)\lambda^{(Q-2)/2}\) in the asymptotic expansion of \(N(\lambda)\)? In other words, is
\[
N(\lambda) = C_{\mathcal H}\lambda^{Q/2} + c_2(\mathcal H)\lambda^{(Q-2)/2} + o(\lambda^{(Q-2)/2})
\]
valid for generic Hörmander sub‑Laplacians?

2. (*Universality*) If such a second term exists, is \(c_2(\mathcal H)\) **universal**—that is, locally computable purely from the nilpotentization and low‑order curvature tensors of the filtered structure (such as those arising in the Tanaka–Webster or general parabolic connections)—and therefore invariant under smooth global deformations preserving the local graded Lie algebra type?  
Or can two sub‑Riemannian manifolds of identical nilpotent type yield distinct \(c_2(\mathcal H)\) values?

Why is it a "good" problem:  
This formulation isolates two profound and unresolved analytical questions at the frontier of hypoelliptic spectral geometry.  
The first—the *existence* of a genuine second‑order spectral term—tests whether the delicate machinery of microlocal analysis (Heisenberg calculus, Rockland parametrices, heat kernel expansions) extends beyond the leading Weyl regime to produce finer asymptotics in the non‑elliptic world.  
The second—the *universality* question—asks if higher‑order spectral invariants are governed purely by the nilpotent tangent data or by deeper geometric curvature features of the distribution.  

Answering either part would decisively advance our understanding of how sub‑Riemannian curvature enters spectral geometry. A positive result would suggest a powerful invariant theory paralleling the Riemannian heat‑kernel coefficients; a negative result would expose previously hidden moduli distinguishing hypoelliptic geometries.  
Its statement remains simple and natural, yet its resolution would require uniting microlocal analysis, representation theory of nilpotent Lie groups, and sub‑Riemannian differential geometry—making it an archetypal **“good” research problem**.